% diversity_reduction_geometric.tex
% Geometric (global) version of the diversity-reduction argument.
% Intended to be \input{} into a main paper, or to replace / supplement
% the Taylor-expansion section in diversity_reduction_theory.tex.
%
% === NOTES ===
%
% This version avoids the first-order Taylor approximation entirely.
% The bounds are exact (not approximate) and hold for arbitrary input
% distributions, not just those with small variance relative to the mean.
%
% The core tool is a chordal-distance contraction lemma for normalization,
% which is elementary but (to our knowledge) has not been applied to
% analyze activation steering before.
%
% The key distinction from the Taylor version: we bound the ABSOLUTE
% post-steering spread, not the CHANGE from pre-steering spread.
% Individual pairs of points can have their angular distance increased
% by steering (see Remark below), but the maximum possible spread after
% steering is bounded and shrinks as O(1/||s||).
%
% The Taylor version from diversity_reduction_theory.tex provides a
% complementary local analysis: it gives the exact covariance structure
% (including which eigenvalues shrink and by how much) but only to first
% order. The geometric version gives global bounds but less structural
% detail. Both are worth including.
%
% Citation keys used (in addition to those in diversity_reduction_theory.tex):
%   \citet{mardia2000directional} -- Mardia & Jupp, "Directional Statistics",
%                                    Wiley, 2000 (for spherical variance)
%
% === END NOTES ===

\subsection{A Global Geometric Bound}
\label{sec:geometric-bound}

The first-order analysis of Section~\ref{sec:covariance} characterizes the
output covariance structure but relies on a local linear approximation.
We now give a complementary global result that holds exactly for arbitrary
input distributions, using only elementary properties of the normalization map.

\subsubsection{Chordal Contraction Under Normalization}

The normalization map $\phi(v) = v/\|v\|$ projects $\mathbb{R}^d \setminus \{0\}$
onto the unit sphere $S^{d-1}$.
The following lemma bounds how far apart two points can land on the sphere
as a function of their Euclidean separation and magnitudes.

\begin{lemma}[Chordal contraction]
\label{lem:chordal}
For any nonzero $u, v \in \mathbb{R}^d$:
%
\begin{equation}
  \|\hat{u} - \hat{v}\|
  \;\leq\;
  \frac{2\,\|u - v\|}{\max(\|u\|, \|v\|)}.
  \label{eq:chordal-bound}
\end{equation}
\end{lemma}

\begin{proof}
Without loss of generality, assume $\|u\| \geq \|v\|$
(the bound is symmetric under swapping $u$ and $v$).
Decompose:
%
\begin{align}
  \hat{u} - \hat{v}
  &= \frac{u}{\|u\|} - \frac{v}{\|v\|}
  \;=\; \frac{u - v}{\|u\|} + v\!\left(\frac{1}{\|u\|} - \frac{1}{\|v\|}\right)
  \label{eq:chordal-decomp}\\
  &= \frac{u - v}{\|u\|} + \frac{v\,(\|v\| - \|u\|)}{\|u\|\,\|v\|}.
  \label{eq:chordal-decomp2}
\end{align}
%
Taking norms and applying the triangle inequality:
%
\begin{equation}
  \|\hat{u} - \hat{v}\|
  \;\leq\;
  \frac{\|u - v\|}{\|u\|}
  + \frac{\|v\| \cdot \big|\|v\| - \|u\|\big|}{\|u\|\,\|v\|}
  \;=\;
  \frac{\|u - v\|}{\|u\|}
  + \frac{\big|\|v\| - \|u\|\big|}{\|u\|}.
  \label{eq:chordal-triangle}
\end{equation}
%
By the reverse triangle inequality, $\big|\|v\| - \|u\|\big| \leq \|u - v\|$, so:
%
\begin{equation}
  \|\hat{u} - \hat{v}\|
  \;\leq\;
  \frac{2\,\|u - v\|}{\|u\|}
  \;=\;
  \frac{2\,\|u - v\|}{\max(\|u\|, \|v\|)}.
  \label{eq:chordal-final}
\end{equation}
\end{proof}

\subsubsection{Concentration Around the Steering Pole}

Let $\phi_s(x) = \widehat{x + s}$ denote the normalized steered representation.
We first bound how far any single point can land from the ``pole'' $\hat{s}$.

\begin{proposition}[Concentration around $\hat{s}$]
\label{prop:concentration}
For any $x \in \mathbb{R}^d$ and any nonzero $s \in \mathbb{R}^d$:
%
\begin{equation}
  \|\phi_s(x) - \hat{s}\|
  \;\leq\;
  \frac{2\,\|x\|}{\|s\|}.
  \label{eq:concentration}
\end{equation}
\end{proposition}

\begin{proof}
Apply Lemma~\ref{lem:chordal} with $u = x + s$ and $v = s$:
%
\begin{equation}
  \|\phi_s(x) - \hat{s}\|
  = \|\widehat{x + s} - \hat{s}\|
  \leq \frac{2\,\|(x + s) - s\|}{\max(\|x + s\|, \|s\|)}
  = \frac{2\,\|x\|}{\max(\|x + s\|, \|s\|)}.
\end{equation}
%
Since $\max(\|x + s\|, \|s\|) \geq \|s\|$, the result follows.
\end{proof}

\noindent
This is the formalization of the geometric intuition: every point on the sphere
is drawn toward the pole $\hat{s}$, with the maximum displacement from the pole
shrinking as $O(1/\|s\|)$.
Critically, this bound is \emph{exact}---it involves no Taylor expansion or
distributional assumptions.

\begin{remark}[The anti-pole]
\label{rem:anti-pole}
If $x \approx -s$, then $x + s \approx 0$ and $\phi_s(x)$ is ill-defined or
can point in an arbitrary direction.
However, for any distribution $\mathcal{D}$ with bounded support
$\|x\| \leq R < \|s\|$, this regime is inaccessible:
$\|x + s\| \geq \|s\| - R > 0$ for all $x$ in the support.
A subtler issue arises when the \emph{mean} satisfies $\mu \approx -s$,
so that the shifted distribution straddles the origin even though no
individual point reaches it.
In this regime, normalization amplifies rather than compresses variation.
This case is analyzed in Section~\ref{sec:conditions}.
\end{remark}

\subsubsection{Pairwise Contraction}

\begin{proposition}[Pairwise chordal bound]
\label{prop:pairwise}
For any $x_1, x_2 \in \mathbb{R}^d$ and nonzero $s$ with
$\|s\| > R \coloneqq \max(\|x_1\|, \|x_2\|)$:
%
\begin{equation}
  \|\phi_s(x_1) - \phi_s(x_2)\|
  \;\leq\;
  \frac{2\,\|x_1 - x_2\|}{\|s\| - R}.
  \label{eq:pairwise}
\end{equation}
\end{proposition}

\begin{proof}
Apply Lemma~\ref{lem:chordal} with $u = x_1 + s$ and $v = x_2 + s$.
Since $(x_1 + s) - (x_2 + s) = x_1 - x_2$, the steering vector cancels
in the numerator:
%
\begin{equation}
  \|\phi_s(x_1) - \phi_s(x_2)\|
  \;\leq\;
  \frac{2\,\|x_1 - x_2\|}{\max(\|x_1 + s\|, \|x_2 + s\|)}.
\end{equation}
%
For the denominator: by the triangle inequality,
$\|x_i + s\| \geq \|s\| - \|x_i\| \geq \|s\| - R$ for each $i$,
so $\max(\|x_1 + s\|, \|x_2 + s\|) \geq \|s\| - R$.
\end{proof}

\noindent
The numerator $\|x_1 - x_2\|$ is the Euclidean distance between the
\emph{unsteered} inputs---the steering cancels.
The denominator grows linearly with $\|s\|$.
Thus, the map $\phi_s$ is Lipschitz on the support of the distribution
with constant $L = 2/(\|s\| - R)$, which becomes a strict contraction
($L < 1$) when $\|s\| > R + 2$.

\begin{remark}[Steering can increase individual pairwise angles]
\label{rem:non-monotone}
The bound \eqref{eq:pairwise} constrains the \emph{absolute} post-steering
angular spread, but does not claim that steering reduces the angular distance
for every pair relative to the unsteered baseline.
A counterexample: in $\mathbb{R}^2$, take $x_1 = (1, 0)$, $x_2 = (0, 1)$,
and $s = (-\tfrac{1}{2}, 0)$.
The unsteered angle between $\hat{x}_1$ and $\hat{x}_2$ is $\pi/2$, but the
post-steering angle is $\arccos(-1/\sqrt{5}) \approx 2.03$ radians---an increase.
This occurs because $s$ moves $x_1$ toward the origin, amplifying angular
sensitivity in that region.
The global bound still holds: as $\|s\|$ grows, such effects vanish because
all points are pulled away from the origin toward the pole.
\end{remark}

\subsubsection{Distribution-Level Diversity Bounds}

Let $\mathcal{D}$ be a distribution over $\mathbb{R}^d$ with bounded support
$\|x\| \leq R$ and let $\phi_s(\mathcal{D})$ denote the pushforward
distribution on $S^{d-1}$.

\begin{corollary}[Spherical variance bound]
\label{cor:spherical-variance}
The spherical variance
$V(\phi_s(\mathcal{D})) \coloneqq 1 - \|\bar{R}\|$
where $\bar{R} = \mathbb{E}_{x \sim \mathcal{D}}[\phi_s(x)]$
is the mean resultant vector \citep{mardia2000directional}, satisfies:
%
\begin{equation}
  V(\phi_s(\mathcal{D}))
  \;\leq\;
  \frac{2\,\mathbb{E}[\|x\|]}{\|s\|}.
  \label{eq:spherical-variance-bound}
\end{equation}
\end{corollary}

\begin{proof}
By Proposition~\ref{prop:concentration}, $\|\phi_s(x) - \hat{s}\| \leq 2\|x\|/\|s\|$ for each $x$.
Taking expectations:
%
\begin{equation}
  \|\bar{R} - \hat{s}\|
  = \|\mathbb{E}[\phi_s(x)] - \hat{s}\|
  = \|\mathbb{E}[\phi_s(x) - \hat{s}]\|
  \leq \mathbb{E}[\|\phi_s(x) - \hat{s}\|]
  \leq \frac{2\,\mathbb{E}[\|x\|]}{\|s\|}.
\end{equation}
%
Since $\|\hat{s}\| = 1$, the reverse triangle inequality gives
$\|\bar{R}\| \geq 1 - \|\bar{R} - \hat{s}\|$, so:
%
\begin{equation}
  V = 1 - \|\bar{R}\| \leq \|\bar{R} - \hat{s}\| \leq \frac{2\,\mathbb{E}[\|x\|]}{\|s\|}.
\end{equation}
\end{proof}

\begin{corollary}[Expected pairwise chordal distance]
\label{cor:pairwise-expected}
For independent $x_1, x_2 \sim \mathcal{D}$:
%
\begin{equation}
  \mathbb{E}\!\left[\|\phi_s(x_1) - \phi_s(x_2)\|\right]
  \;\leq\;
  \frac{4\,\mathbb{E}[\|x\|]}{\|s\|}.
  \label{eq:pairwise-expected}
\end{equation}
\end{corollary}

\begin{proof}
By the triangle inequality through the pole:
$\|\phi_s(x_1) - \phi_s(x_2)\| \leq \|\phi_s(x_1) - \hat{s}\| + \|\hat{s} - \phi_s(x_2)\|$.
Applying Proposition~\ref{prop:concentration} to each term and taking expectations
over independent $x_1, x_2$:
%
\begin{equation}
  \mathbb{E}\!\left[\|\phi_s(x_1) - \phi_s(x_2)\|\right]
  \leq \frac{2\,\mathbb{E}[\|x_1\|]}{\|s\|} + \frac{2\,\mathbb{E}[\|x_2\|]}{\|s\|}
  = \frac{4\,\mathbb{E}[\|x\|]}{\|s\|}.
\end{equation}
\end{proof}

\begin{corollary}[Spherical diameter]
\label{cor:diameter}
The diameter of $\phi_s(\mathcal{D})$ in chordal distance satisfies:
%
\begin{equation}
  \mathrm{diam}\!\left(\phi_s(\mathcal{D})\right)
  \;\leq\;
  \frac{4R}{\|s\|},
  \label{eq:diameter}
\end{equation}
%
where $R = \sup\{\|x\| : x \in \mathrm{supp}(\mathcal{D})\}$.
% NOTE: This uses the triangle-inequality-through-the-pole route (hence 4R
% instead of 2D_eucl). The direct Proposition 3 bound gives the tighter
% 2 * diam(D) / (||s|| - R), but requires ||s|| > R. For large ||s||
% relative to R, both are O(R/||s||).
\end{corollary}

\begin{proof}
For any $x_1, x_2$ in the support, Proposition~\ref{prop:concentration} gives
$\|\phi_s(x_i) - \hat{s}\| \leq 2R/\|s\|$.
By the triangle inequality,
$\|\phi_s(x_1) - \phi_s(x_2)\| \leq 4R/\|s\|$.
\end{proof}

\subsubsection{Interpretation}

These results formalize the geometric picture of diversity reduction
under steering.
The distribution of post-normalization representations is confined to a
spherical cap of chordal radius $2R/\|s\|$ centered at the pole $\hat{s}$.
As the steering magnitude increases, this cap shrinks, and all measures
of angular dispersion---spherical variance, expected pairwise distance,
diameter---decrease as $O(1/\|s\|)$.

\paragraph{Comparison with the local analysis.}
The first-order expansion of Section~\ref{sec:covariance} and the global
bounds of this section are complementary:

\begin{center}
\begin{tabular}{lcc}
  \toprule
  & \textbf{Local (Taylor)} & \textbf{Global (geometric)} \\
  \midrule
  Approximation & First-order; requires $\|\delta\| \ll \|m\|$ & Exact \\
  Diversity measure & Effective rank, eigenvalues of $\Sigma_z$ & Spherical variance, chordal distance \\
  Structural detail & Full covariance structure & Scalar bounds only \\
  Applicable regime & Large $\|s\|$ (increasingly accurate) & All $\|s\| > 0$ \\
  \bottomrule
\end{tabular}
\end{center}

\noindent
The local analysis reveals \emph{which} dimensions of variation are lost
(those along $\hat{m}$) and by how much (the eigenvalues of
$P_{\perp m}\Sigma_x P_{\perp m}/\|m\|^2$).
The global analysis establishes that post-normalization spread is \emph{bounded}
for any steering magnitude, without approximation error, and gives
quantitative scaling with $\|s\|$.
Whether this bound implies a \emph{reduction} relative to the unsteered
baseline depends on additional conditions analyzed in
Section~\ref{sec:conditions}.